\documentclass[aspectratio=169,12pt]{beamer}
\usepackage{amsmath}
\usepackage{amssymb}
\usepackage{array}
\usepackage[most]{tcolorbox}
\usepackage{graphicx}
\usepackage[sfdefault,lf]{FiraSans}
\renewcommand{\familydefault}{\sfdefault}
\setlength{\parindent}{0pt}
\everymath{\displaystyle}

\setbeamertemplate{navigation symbols}{}
\setbeamertemplate{footline}{}
\setbeamertemplate{headline}{}
\setbeamertemplate{blocks}[rounded][shadow=false]

\definecolor{mmpageWhite}{HTML}{FCFBF7}
\definecolor{mmtanSoft}{HTML}{F7F5EF}
\definecolor{mmgreenDeep}{HTML}{244B2C}
\definecolor{mmgreenLeaf}{HTML}{5D8A56}
\definecolor{mmgreenSoft}{HTML}{E4EFD9}
\definecolor{mmtext}{HTML}{163221}
\definecolor{mmwarn}{HTML}{8F2F36}
\definecolor{mmwarnSoft}{HTML}{F8E8EA}

\setbeamercolor{background canvas}{bg=mmpageWhite}
\setbeamercolor{normal text}{fg=mmtext,bg=mmpageWhite}
\setbeamercolor{block title}{fg=mmgreenDeep,bg=mmgreenSoft}
\setbeamercolor{block body}{fg=mmtext,bg=mmtanSoft}
\setbeamercolor{block title example}{fg=mmgreenDeep,bg=mmgreenSoft}
\setbeamercolor{block body example}{fg=mmtext,bg=white}
\setbeamercolor{block title alerted}{fg=mmwarn,bg=mmwarnSoft}
\setbeamercolor{block body alerted}{fg=mmtext,bg=white}
\setbeamercolor{itemize item}{fg=mmgreenDeep}
\setbeamercolor{itemize subitem}{fg=mmgreenDeep}

\newtcolorbox{MMCard}[2][]{enhanced,breakable=false,colback=#2,colframe=black,boxrule=1.5pt,arc=8pt,left=5pt,right=5pt,top=4pt,bottom=4pt,before skip=0pt,after skip=0pt,#1}
\newcommand{\KoruIcon}{\raisebox{-0.2em}{\includegraphics[height=1.05em]{../../../manamaths/SITE/header-logo.png}}}
\newcommand{\NotesTitle}[1]{{\colorbox{mmtanSoft}{\parbox{0.975\linewidth}{\vspace{0.14em}\hspace{0.25em}{\LARGE\bfseries\textcolor{mmgreenDeep}{#1}\hfill\KoruIcon}\vspace{0.14em}}}\par\vspace{0.18em}{\color{mmgreenLeaf}\rule{\linewidth}{2.0pt}}\par\vspace{0.12em}}}
\newcommand{\SectionTitle}[1]{{\large\bfseries\textcolor{mmgreenDeep}{#1}}\par\vspace{0.04em}}
\newcommand{\GreenDivider}{\par\vspace{0.01em}{\color{mmgreenLeaf}\rule{\linewidth}{2.4pt}}\par\vspace{0.03em}}
\newcommand{\KeyIdea}[1]{\begin{MMCard}[title={\textbf{Key idea}},colbacktitle=mmgreenSoft,coltitle=mmgreenDeep]{mmtanSoft}\small #1\end{MMCard}}
\newcommand{\WorkedExample}[2]{\begin{MMCard}[title={\textbf{#1}},colbacktitle=mmgreenSoft,coltitle=mmgreenDeep]{white}\small #2\end{MMCard}}
\newcommand{\CommonMistake}[1]{\begin{MMCard}[title={\textbf{Common mistake}},colbacktitle=mmwarnSoft,coltitle=mmwarn]{white}\small #1\end{MMCard}}
\newcommand{\TryThese}[1]{\begin{MMCard}[title={\textbf{Try these}},colbacktitle=mmgreenSoft,coltitle=mmgreenDeep]{mmtanSoft}\small #1\end{MMCard}}
\newcommand{\StepLine}[2]{\textbf{#1.} #2\par\vspace{0.03em}}

\usepackage{tikz}

\setbeamertemplate{background}{
 \begin{tikzpicture}[remember picture,overlay]
 \draw[
 black,
 line width=1.5pt,
 rounded corners=10pt
 ]
 ([xshift=0.45cm,yshift=-0.45cm]current page.north west)
 rectangle
 ([xshift=-0.45cm,yshift=0.45cm]current page.south east);
 \end{tikzpicture}
}

\begin{document}

\begin{frame}[t,shrink=8]
\NotesTitle{Notes \& Steps}
\footnotesize
\begin{columns}[T,onlytextwidth]
\begin{column}{0.48\textwidth}
\KeyIdea{When you multiply or divide by 10, 100, or 1000, the digits move left or right in their place-value columns. Multiply shifts left; divide shifts right.}
\SectionTitle{Steps for multiplying}
\StepLine{1}{Count the zeros in 10, 100, or 1000.}
\StepLine{2}{Move the digits that many places to the left.}
\StepLine{3}{Fill any empty columns with zeros.}
\SectionTitle{Steps for dividing}
\StepLine{1}{Count the zeros.}
\StepLine{2}{Move the digits that many places to the right.}
\StepLine{3}{For decimals past the point, keep the decimal point fixed.}
\end{column}
\begin{column}{0.48\textwidth}
\SectionTitle{Pattern examples}
\begin{itemize}
  \setlength{\itemsep}{0.12em}
  \item \(6 \times 10 = 60\) (6 moves left one place)
  \item \(4 \times 100 = 400\) (4 moves left two places)
  \item \(80 \div 10 = 8\) (8 moves right one place)
  \item \(500 \div 100 = 5\) (5 moves right two places)
  \item \(3.4 \times 10 = 34\) (decimal moves right)
  \item \(5.8 \div 10 = 0.58\) (decimal moves left)
\end{itemize}
\end{column}
\end{columns}
\GreenDivider
\vfill
\begin{columns}[b,onlytextwidth]
\begin{column}{0.48\textwidth}
\CommonMistake{"Just add a zero" only works for whole numbers. For example, \(3.4 \times 10 \neq 3.40\). The answer is \(34\), not \(3.40\). The digits shift, zeros are not just tacked on.}
\end{column}
\begin{column}{0.48\textwidth}
\TryThese{\begin{enumerate}
  \item \(9 \times 1000\)
  \item \(7000 \div 1000\)
  \item \(0.7 \times 100\)
\end{enumerate}}
\end{column}
\end{columns}
\end{frame}

\begin{frame}[t,shrink=12]
\NotesTitle{Notes \& Steps}
\begin{columns}[T,onlytextwidth]
\begin{column}{0.48\textwidth}
\WorkedExample{Example 1: multiplying decimals}{Calculate \(1.25 \times 10\).\[1.25\]Move the digits one place left: \(12.5\). Answer: \(12.5\).}
\end{column}
\begin{column}{0.48\textwidth}
\WorkedExample{Example 2: dividing decimals}{Calculate \(4.2 \div 100\).\[4.2\]Move the digits two places right: \(0.042\). Answer: \(0.042\).}
\end{column}
\end{columns}
\GreenDivider
\TryThese{\begin{enumerate}
  \item \(3.6 \times 100\)
  \item \(0.48 \times 1000\)
  \item \(72 \div 100\)
\end{enumerate}}
\end{frame}

\begin{frame}[t,shrink=12]
\NotesTitle{Notes \& Steps}
\begin{columns}[T,onlytextwidth]
\begin{column}{0.48\textwidth}
\WorkedExample{Example 3: dividing a small decimal}{Calculate \(0.09 \times 1000\).\[0.09\]Move three places left. \(9\) needs zeros: \(09.0 \to 90\). Answer: \(90\).}
\end{column}
\begin{column}{0.48\textwidth}
\WorkedExample{Example 4: applying to real units}{Convert \(0.047\) km to metres.\[0.047 \times 1000 = 47\]So \(0.047\) km = \(47\) m.}
\end{column}
\end{columns}
\GreenDivider
\CommonMistake{When dividing, digits move right — the number gets smaller. Many students think \(4.2 \div 100 = 420\) because they add zeros. Actually \(4.2 \div 100 = 0.042\).}
\end{frame}

\end{document}
